% This must be in the first 5 lines to tell arXiv to use pdfLaTeX, which is strongly recommended.
\pdfoutput=1
% In particular, the hyperref package requires pdfLaTeX in order to break URLs across lines.

\documentclass[11pt]{article}

% Remove the "review" option to generate the final version.
\usepackage[review]{emnlp2021}

% Standard package includes
\usepackage{times}
\usepackage{latexsym}
\usepackage{todonotes}

% For proper rendering and hyphenation of words containing Latin characters (including in bib files)
\usepackage[T1]{fontenc}
% For Vietnamese characters
% \usepackage[T5]{fontenc}
% See https://www.latex-project.org/help/documentation/encguide.pdf for other character sets

% This assumes your files are encoded as UTF8
\usepackage[utf8]{inputenc}

% This is not strictly necessary, and may be commented out,
% but it will improve the layout of the manuscript,
% and will typically save some space.
\usepackage{microtype}
\usepackage{graphicx}
\usepackage{import}
\usepackage{layout}
\usepackage{tabularx, makecell}
\usepackage{booktabs}
\usepackage{mathrsfs}
\usepackage{amssymb} 
\usepackage{url}
\usepackage{graphicx}
\usepackage{xspace,paralist}
\usepackage{times,latexsym}
\usepackage{amsmath}
\usepackage{appendix}
\usepackage{comment} 
\usepackage{enumitem}
\usepackage{makecell}
\usepackage{multirow}
\usepackage{xcolor}
\usepackage{cleveref}
\newcommand{\eqnref}[1]{Eq~\eqref{#1}\xspace}
\newcommand{\tabref}[2][]{Table#1~\ref{#2}\xspace}
\newcommand{\figref}[1]{Figure~\ref{#1}\xspace}
\newcommand{\secref}[1]{Section~\ref{#1}\xspace}
\newcommand{\appref}[1]{Appendix~\ref{#1}\xspace}
% If the title and author information does not fit in the area allocated, uncomment the following
%
%\setlength\titlebox{<dim>}
%
% and set <dim> to something 5cm or larger.

\title{Factcheck-GPT: A Holistic Framework for Fact-Checking\\ the Output of Generative Large Language Models}
% \title{Factuality in Large Language Model Generation:\\
% An Annotation Framework and Challenges}

% Author information can be set in various styles:
% For several authors from the same institution:
% \author{Author 1 \and ... \and Author n \\
%         Address line \\ ... \\ Address line}
% if the names do not fit well on one line use
%         Author 1 \\ {\bf Author 2} \\ ... \\ {\bf Author n} \\
% For authors from different institutions:
% \author{Author 1 \\ Address line \\  ... \\ Address line
%         \And  ... \And
%         Author n \\ Address line \\ ... \\ Address line}
% To start a seperate ``row'' of authors use \AND, as in
% \author{Author 1 \\ Address line \\  ... \\ Address line
%         \AND
%         Author 2 \\ Address line \\ ... \\ Address line \And
%         Author 3 \\ Address line \\ ... \\ Address line}

\author{First Author \\
  Affiliation / Address line 1 \\
  Affiliation / Address line 2 \\
  Affiliation / Address line 3 \\
  \texttt{email@domain} \\\And
  Second Author \\
  Affiliation / Address line 1 \\
  Affiliation / Address line 2 \\
  Affiliation / Address line 3 \\
  \texttt{email@domain} \\}

\begin{document}
\maketitle
\begin{abstract}
The increased use of large language models (LLMs) across a variety of real-world applications calls for mechanisms to verify the factual accuracy of their outputs. In this paper, we present a comprehensive framework for annotating the factuality of LLM-generated responses. Our framework encompasses a multi-stage annotation scheme that is designed to yield detailed labels concerning the verifiability and factual inconsistencies found in LLM outputs. The annotation procedure delineates associated tasks for detecting and correcting factual errors in LLM generations, such as decontextualization, checkworthiness identification, atomic claim decomposition, external knowledge verification, and fine-grained error correction. We employ our framework in an annotation study that highlights and examines several critical challenges such as distinguishing objective facts from subjective opinions, navigating ambiguity in identifying fact hallucinations, and making correctness judgements in the presence of conflicting evidence.
Our work emphasizes the intricate nature of fact-checking machine-generated language within the context of LLMs and offers insights for future development of more accurate and reliable LLMs.

%This study aims to frame annotation guidelines and to clarify related tasks for automatically detecting and correcting factual errors generated by LLMs such as ChatGPT.

\end{abstract}

\section{Motivation}
Large language models (LLMs) have demonstrated impressive capabilities in terms of generating naturally sounding text. 
Particularly ChatGPT, through its scale and reinforcement learning with human feedback (RLHF), has garnered significant attention.
Its dialog interface allows users to interact with the underlying LLM more effectively and efficiently via interactive chats.
Fluent and comprehensive answers over a broad range of human inquiries surpass prior public chatbots in terms of usefulness.
However, despite its powerful capabilities, several studies have consistently shown significant remaining challenges, including errors in terms of factuality, reasoning, math, coding, biases, etc.~\citep{bang2023multichatgpt, ali2023failurechatgpt, giuven2023llmsfailures}.
We focus on factual errors and aim to automatically detect and correct these errors in the generated text, with an emphasis on long-document generations, thus enhancing the factuality of future LLMs.


Existing fact-checking datasets either focus on factoid of human-written text or attend to hallucination of sentence-level generations by LLMs.
To close this gap of analysing factual errors of document-level machine-generated text by LLMs, we construct the first document-level fact-checking corpus of LLMs, over which we conduct an in-depth analysis.
% Baseline model based on previous fact-checking methods shows large headroom for improvement.


\section{Scope of Generations}
The first question to consider is \emph{what kind of generations are we most concerned about?}

\paragraph{Subjectivity} While we can ask a lot of questions ranging across diverse domains and we can analyse the corresponding answers generated by the LLMs, answers including more objective factual statements are more interesting than subjective opinions: it makes very little sense to attempt to check the factual correctness of a personal opinion.
Moreover, objective claims such as that \emph{the sun rises from the East} are also not what we care about.
We are primarily interested in texts where LLMs are prone to hallucinate.   
To this end, the goal of this work is to collect document-level generations, where the majority of the statements are making objective statements and are thus prone hallucinations. 

The next questions we can ask are: \textit{1. Which domain questions tend to result in objective answers?} and \textit{2. When LLMs tend to hallucinate?}

Unfortunately, we do not have a comprehensive answer. 
We can only glimpse a tip of the iceberg by some probing studies.
For the first question, news, Wikipedia and question-answering websites that provide professional knowledge such as Stack Overflow for programming questions might be more factoid-intensive than poetry, lyrics, and stories. 
For the latter, LLMs are shown to have limited memorisation for less frequent entities, and are prone to hallucinate on long-tail knowledge~\citep{nikhil2022longtail, alex2022when}.
Therefore, we might be able to ask questions involving long-tail knowledge.

The following step is to decide \textit{Where and how to collect data?}\todo{We need to address this question.}


\section{Expected Output}
Imagine we ask ChatGPT a question Q, and it gives an answer A: what our concern about this answer could be?
% Given a pair of (Q, A), where Q is a question, A is a generated answer by a LLM, what is our concern of the answer?
We do not know \textit{when to trust language models and when not to}.
Thus, we can expect taht there will be a reliable system that can tell us \textbf{how much we can trust each piece of information} provided by ChatGPT and to offer evidence to support that trustworthiness value.
Meanwhile, when LLM-generated statements are non-factual, the system can correct them.

\paragraph{How much vs. Yes-or-No}
That is, instead of detecting whether a claim is factual or not by yes-and-no answer, we expect the system to be uncertainty-aware. 
For example, given a claim, there are 100 pieces of retrieved evidence, 80 among them support the claim, 10 refute it, and 10 just discuss it.
We will have an empirical probability distribution (0.8. 0.1, 0.1), indicating the latent correctness likelihood of this claim, without distinguishing the reliability weight of each evidence source. 
Put differently, there is 80\% confidence to affirm this claim can be trusted, depending on open-domain world knowledge.

\paragraph{When to correct?}
After identifying the trustworthiness probability of a claim, the next question is \emph{When to correct?}

The recently proposed RARR model~\citep{gao2022attributed}, decides when to correct an input, i.e., the entire answer rather than a claim, using an agreement model.
Specifically, given a query that may cover some aspects of an answer, if the evidence and the partially-revised input give the same answer, then they consider that the evidence agrees with the current statement; otherwise, they consider that there is a factual error, and they propose an edit.
Note that they implicitly assume that each piece of evidence is factually correct and that there are no contradictions between the different pieces of evidence. However, contradicting stance are quite frequent in practice.

One option is to use the confidence distribution and to propose to activate the correction module when the probability of the supporting evidence is less than a threshold, e.g., $\tau<$ 0.7.
We may need a case study to investigate how to set an appropriate threshold.


% Highlight high-level novelty of the framework.
\subsection{Why is uncertainty-aware output better?}
Check-worthy claims or claims where people are not sure and require affirmation, are mostly boundary and ambiguous cases.
They are either difficult for people to judge as true or false, i.e.,~requiring expert-level knowledge, or they involve controversial topics that are hard to reach consensus on depending on the available evidence.
Framing fact-checking labels into uncertainty-aware distribution (soft labels) instead of hard ones, fits the uncertain nature of fact-checking.

Moreover, akin to LLMs' parametric knowledge (i.e.,~the knowledge stored in their parameters), can we completely trust non-parametric knowledge (i.e., the retrieved text chunks)?
Soft labels on the other hand give high tolerance of the quality of the collected evidence, as any evidence does not play the dominant role determining the final label, but partial elements supporting the final decision.
They do not neet to be totally reliable or of high quality. 


\section{Problem Definition}
We propose the task of \textbf{uncertainty-aware} fact-checking \textbf{document-level} generations of LLMs, including checking the factuality of long documents, correcting factual errors, and providing evidence in teh form of citations.
This differs from previous work mainly in two perspectives.

% Challenges and novelty of this task
\paragraph{Document Level}
Compared to sentence-level claims, for which we can easily estimate check-worthiness, directly retrieve relevant evidence, and make a decision, a major challenge for document-level fact-checking is how to decompose the text into individual factual claims and then how to decontextualise the statements.

\textit{Our setup vs. RARR}
Skipping decomposition and decontextulisation, RARR achieves fact-checking over long documents in two steps: generating questions covering different aspects of the input, and retrieving evidence for each of the questions it has generated.
They collects evidence from the entire document.
Both the detection and the correctness are performed from the viewpoint of the entire answer, rather than of a fine-grained claim.
It effectively avoids the challenges we mentioned above.

However, firstly it may miss some errors when generated questions do not cover all aspects of an answer, and such risk will be increased with documents becoming longer.
Secondly, detection and revision over whole answer fail to locate factual errors and associate related evidence to a specific claim. 
We don't know which claim contains factual errors and what makes it non-factual.
Additionally, it is prone to revise too much.
It cannot locate to a specific error, so the model does not know where and to what extent the original text should be revised.
This leads RARR to use preservation metric to control, so that the revised text preserves aspects of the original text.

\paragraph{Uncertainty-Awareness}
Different from previous work assuming that all retrieved evidence are reliable and correct, acknowledge different stances of evidence in terms of a given claim.
We frame the label of fact checking as a soft probability that illustrates the distribution of how real-world knowledge supports/refutes a claim, instead of a hard label indicating whether a claim itself is factual or not.

For example, 10 evidence are collected after retrieval given a claim, where eight support the claim, one refutes and one stands neutral.
Previous work are prone to apply the method of majority voting and present claim is factual.
However, this aggregation process loses information.
The final label neither reflects full information that it may bias the latent truth, nor matches the uncertain nature of diverse stances of real evidence.
We will show empirical distribution of (0.8, 0.1, 0.1) for this claim, demonstrating that the majority evidence support it is true while there is also opposite opinions. 
The final decision gives to users rather than the automatic system, as the degree of factual requirements differs among users and applications.

\paragraph{Nature of the Edits} We differ from RARR as we allow multiple edits in a claim (e.g., fixing simultaneously the conference name, the year, and the nature of the award, e.g., for Preslav's best demo award). Moreover, we allow dropping a claim entirely, e.g., when the claims is that somebody was an ACL Fellow, but this never happened: so, instead of saying that the person did NOT get this recognition, we just drop it entirely.

\section{Pipeline of Data Collection and Subtasks}
\figref{fig:pipeline} presents the overview of the task, coupled with an example showing outputs for each step.
\begin{figure*}[t!]
	\centering
	\includegraphics[scale=0.4]{PipelinewithExample.pdf}
	\caption{\textbf{Left:} Overall process of fact-checking a long document generated by ChatGPT. \textbf{Right:} An example flow over the pipeline.}
	\label{fig:pipeline}
\end{figure*}
\paragraph{Decompose and Decontextualise}

\todo{Add a note that the decomposition can be done by ChatGPT, and then postedited by a human annotator. In this, case, we need to record what the original decomposition was, and what was postedited. This will allow us to evaluate ChatGPT for this task, and to discuss it.}
Given a long-document answer $A$ generated by ChatGPT, it is infeasible to fact-check the whole document at once.
The first step is to decompose the answer into a set of atomic checkable context-independent statements. 
% YXW: What is checkable? What is not?
This poses several challenges. 
(1) How to decompose the answer? 
(2) How to determine the granularity of the decomposition? 
In other words, when and where to decompose?
For example, for the claim \textit{Canada is a constitutional monarchy, and as such, it does not have a king.}, it can be fact-checked as one statement, or be decomposed into two sub\-claims: \textit{Canada is a constitutional monarchy.} and \textit{It does not have a king.} 

Moreover, how to decontextualise or reframe the claim that it is context-independent? 
Statements in the answer might be context-dependent. 
Discourse and co\-reference relations exist between sentences.
For example, it is invalid to check statement \textit{It does not have a king} before replacing ``It'' with ``Canada'' or ``Constitutional monarchy''.   
In addition to co\-reference relation, for the fourth statement, it is not reasonable to check the claim \textit{Queen Elizabeth II is also the queen of 15 other Commonwealth realms.}
Instead, the claim should be re\-framed to \textit{Queen Elizabeth II is the queen of 16 Commonwealth realms (including Canada).}
Afterwards, the next step is to identify whether the statement is check-worthy or not.

\paragraph{Identify check\-worthy claims}
Not all statements in the answer require fact-checking. 
Each statement needs to be identified as being a claim or an opinion.
We have only taken account of objective fact against subjective judgement.
\textbf{Any other aspects to be considered?}
For instance, the role (importance) of the claim to the answer is also a crucial criteria for its check-worthiness.
For example, the second sentence needs more attention than the last sentence.

After omitting statements that contain personal opinions, we gather a set of check-worthy claims.
Then, we need to verify the claim by retrieving and collecting evidence.

\paragraph{Retrieve and Collect evidence/citations}
\textit{How and where to retrieve evidence?}
We can use Google search engine, or only search from wikipedia.
We can choose to search by questions covering different aspects of the claim, alternatively claim itself or entities in the claim.

\todo{Note that we can do the query generation using ChatGPT.}

\begin{comment}
Though incorporating retrieved evidence largely helps address LLMs' hallucinations, it is unclear whether it is strictly superior or complementary to parametric knowledge~\citep{liu-etal-2022-token}.
We may be able to leverage both: non-parametric knowledge by Google search engine and parametric knowledge via properly prompting ChatGPT.
Given a check-worthy claim, we can collect web pages with citations as well as knowledge which has been memorised by LLMs.
\end{comment}

\paragraph{How well does ChatGPT perform over aforementioned subtasks?}
ChatGPT shows potential to decompose a long answer into atomic statements, and further identify whether they are objective facts that can be verified or facts that contain personal opinions.
However, there is no guarantee that separated claims are context-independent and the correctness of the split of objective vs.\ subjective statements.
So outputs of ChatGPT can only be regarded as preliminary results, being provided to annotators to validate and confirm. 

For collecting evidence with related URL links supporting a claim, not surprisingly, all URL links generated by ChatGPT are wrong.
This also implies the importance of incorporating external knowledge with citations to support claims.


\paragraph{Identify stance of evidence}
How to identify the stance of evidence? 
RARR achieves this by assessing whether answers depending on the evidence and the claim are same or not, in terms the given query. 
If they are same, then the evidence support the claim, otherwise refutes.
Previous work also employ NLI model to classify whether the claim is entailed by evidence, or controversial against evidence, or not enough information is provided to judge.
What method will be apply? and Why?

After obtaining stance of each piece of evidence, we can statistically calculate uncertainty-aware labels for claims, and further obtain soft factual label for a whole answer.
So here is an open problem: \textit{how to decide soft label of a whole answer depending on elementary labels of component claims?}

Besides, \textbf{do we consider to provide the type of hallucination}, such as lacking domain-specific knowledge, commonsense knowledge, unrelated to the central topic, conflict to the preceding or succeeding context and etc~\citep{liu-etal-2022-token}.


\paragraph{Correct claim and Deduplicate answer}
\textit{What criteria should we rely on} when correcting each non-factual claim depending on related evidence, such as preservation. 
Then, we merge all statements in order either the originally-factual one or the correctly-revised one. 
Finally paraphrasing is needed to delete reduplicative content and recover a fluent answer.


% \paragraph{Validate the revised answer}


\subsection{Questions to Discuss}
\begin{itemize}
    \item What is the scope of LLMs' generations? What domains and questions are we more interested in?
    \item What LLM would we like to include?
    \item How to decompose and decontextualise? What is the standard? What kind of criteria do we rely on in annotation?
    \item How and where to collect evidence?
    \item What criteria will activate the edit model to correct the claim?
    \item What's the standard or guidelines for claim correction? What criteria should we rely on in claim correction?
    \item How to recover a fluent answer and validate its correctness?
\end{itemize}


\section{Annotation Guideline}
For each example, annotators are given a pair of (query, response).
A response is either an answer generated by LLMs responding to users' question, or a document returned by LLMs according to users' request.
Annotators are required to give outputs of each step shown in \figref{fig:pipeline}. 
We describe how to annotate for component subtasks throughout the pipeline, particularly clarify how to deal with possible ambiguous scenarios.

\subsection{Decompose}
It is subjective to decide the granularity of decomposition.
We may aim to break down a long document into a set of atomic claims, while the definition of a atomic claim varies.
Here, we practically apply the following strategy: 
\begin{itemize}
    \item Start by decomposing into sentences.
    \item If the sentence contains too much information, break it into several components, but annotators do not overdo it, e.g., decomposing \textit{Capitol Hill riots happened on January 6, 2021} to one claim for year and one for day. 
    \item If several pieces of information are strongly dependent with each other, they are expected to co-occur in one snippet of evidence text, no more breaking-down is needed.
\end{itemize}


\subsection{Decontextualise}
The criteria of decontextualisation is to ensure that all separated statements fully preserve semantics presented in original context.
For example, statement that \textit{it happened on Jan 6, 2021} loses information in decomposition, which makes it uncheckable. 
In such cases, annotators should replace pronouns, such as \textit{it, they, those, these, this, that}, with specific entities or events after decomposition.
Decontextulisation is mostly needed over cases with co\-reference relation.
For complex relation, such as two sentences are strongly dependent with each other, we encourage to go back to step of decomposition and keep the original text without breaking-down.


\subsection{Identify check-worthy claim}
We consider two aspects in check-worthiness identification:
\begin{itemize}
    \item If a statement presents subjective opinions, then it is not check-worthy. % otherwise move to next step.
    \item If the objective facts presented in a statement are commonsense, such as \textit{sun rises from the east}, it is not worth checking. 
\end{itemize}
Therefore, we regard a statement as check-worthy claim when it presents objective facts, and these facts are not naive commonsense.
There is a special case. If the objective facts presented in a statement are not publicly available information. Namely, we cannot collect any evidence over web pages related to the claim, such as personal experience.
They are regarded as uncheckable claims.

Annotators are asked to identify each statement by four classes, and count the number of statements falling into each class for a response.
\begin{itemize}
    \item[] 1. checkable and checkworthy claim
    \item[] 2. subjective opinion
    \item[] 3. commonsense
    \item[] 4. uncheckable fact
\end{itemize}


\subsection{Retrieve and collect evidence}
Given a claim, based on general web pages (including Wikipedia pages), annotators are asked to record all queries tried in retrieval (e.g. questions covering some aspects of a claim, or entities contained in a claim), and indicate those which are finally used to get helpful evidence.

In evidence collection, annotators need to gather 10 most relevant snippets of text that can verify the factuality of the claim, from either one or multiple web pages.
For each snippet of evidence, record meta-info including (1) query, (2) citation (URL) of the web page from which this piece of evidence is extracted, (3) personal judgement of whether the source of evidence is reliable or not. it's a continuous value in range of [0, 1], where 1 indicates completely reliable and 0 indicates completely unreliable.\footnote{Source reliability can also automatically be collected from MBFC/AllSides/Politifact/, but they apply for a small number of sources.} and (4) indicator whether this individual evidence is sufficient to verify the input claim. 

Note that 10 pieces of evidence is not a hard criteria.
If less than 10 (even only one) pieces of evidence are sufficient to verify the input claim, and they are from reliable sources, annotators are allowed to collect <10 results.
Meanwhile, if a claim involves a controversial topic, annotators are also encouraged to collect more than 10 results.

Aforementioned guidelines are applicable to claims for which there exist evidence over web pages.
However, there are situations where there isn't any information in public web pages, e.g. personal experience.
They are objective facts, but are not extensively known by the public.
Put differently, they are uncheckable.
Here, we suggest stepping back to the last step and categorising as uncheckable claims. 


\subsection{Identify evidence stance}
Given a claim and K evidence results, raters are asked to annotate each evidence by a soft label (s, r, n), representing the percentage of information in the claim is supported, refuted or neutral against the evidence.

% \begin{comment}
\paragraph{Discussion} We have two questions to discuss.
\\
Q: Is stance of \textit{neutral} necessary to set in fact-checking context? When is it unimportant? When does it matter? \\
A: Our goal is to verify whether a claim is factual or not.
Neutral evidence is the information related to the claim, but is not enough to support or refute any information in a claim. If so, it possibly shouldn't be collected as evidence if there exist evidence either supporting or refuting a claim.
But if there isn't any direct evidence, neutral information is also important.
\\
\\
Q: Do we need extra label of \textit{partially support}?\\
A: We may prefer to explicitly keep three-class labelling in annotation, but implicitly achieve this by assigning soft labels. 

There are mainly two reasons. 
Firstly, we have partially support, then we correspondingly have partially refute, too many labels tend to confuse raters. 
The simpler the guideline is, the easier for annotation to label.
Secondly, this issue is essentially attributed to that the claim wasn't broken into fine-grained subclaims. 
Term of ``partially'' is a vague notion.
Supporting 20\%, 50\% or 80\% content of the claim all can be labelled as partially support, but obviously, they differs and influence following determination of factuality degree of a claim.
It may be sensible to give a specific value indicating the percentage of information in a claim is supported or refuted by the evidence. 

For the example that triggered requirement of \textit{partially support}, 
Claim: Preslav Nakov has received the Bulgarian President’s John Atanasoff Award for outstanding young Bulgarian computer scientists (2011).
\begin{itemize}
    \item[-] E1: He was also the first to receive the Bulgarian President's John Atanasoff award, named after the inventor of the first automatic electronic digital computer.
    \item[-] query used to retrieve: Preslav Nakov
    \item[-] source URL: \url{https://mbzuai.ac.ae/study/faculty/preslav-nakov/}
    \item[-] source reliability: Yes (reliable)
    \item[-] individual evidence sufficiency: No (not sufficient to verify the claim)  
\end{itemize}
There are three major subclaims to verify in this claim: award name, reason and year.
The evidence support one element, assigning (0.33, 0, 0).
% \end{comment}

\subsection{Determine factuality label}
\paragraph{Claim label:}
Afterwards, for each claim, we have K evidence results, including a set of snippets of text, corresponding stance vectors $[\mathbf{s}_1, \mathbf{s}_2, \dots, \mathbf{s}_K]$ and source reliability values $[r_1, r_2, \dots, r_K]$.
We record two factuality labels and their corresponding normalised versions for each claim.
One simply sums up stances, and the other incorporates weights of reliability as below:
\begin{itemize}
    \item[(1)] $\mathbf{l_{s}} = \sum_{i=1}^{K} \mathbf{s_i}$
    \item[(2)] $\mathbf{l_{sr}} = \sum_{i=1}^{K} r_i \mathbf{s_i}$
\end{itemize}

\paragraph{Document label:}
Factual labels of the whole document consider four different types of component statements.
\begin{itemize}
    \item check-worthy and checkable statement, we use their normalised factual label $\mathbf{l_{sr}}$; 
    \item commonsense statement, it implies that all available evidence support it, we use $\mathbf{l_{sr}}=(1,0,0)$;
    \item uncheckable facts, there isn't any information supporting or refuting such a statement, all web pages cannot provide enough information and keep neutral, we use $\mathbf{l_{sr}}=(0,0,1)$;
    \item subjective opinions, they do not have the characteristic of factuality, so they do not have factual labels.
\end{itemize}
% Assuming there are n1 checkable and check-worthy claims in the answer, and n2 commonsense statements, and n3 uncheckable facts, and n4 subjective-opinion statements, 
Then the factual label of a document is aggregated as $\mathbf{l_{d}} = \sum \mathbf{l_{sr}}$.


\paragraph{Discussion} 
What kind of normalisation should we use, do different normalisation strategies affect the meaning of the final results significantly?

Should we consider subjective opinions when determine the factual label of the whole response/document?
Why we should or why not?
If we consider, how to incorporate them? 

\subsection{Determine correction}
Based on a set of evidence and the normalised factuality label $\mathbf{l_{sr}}=(s, r, n)$ for a claim, we follow criteria below to decide whether we will perform editing.
\begin{itemize}
    \item If the claim is completely supported by evidence, i.e. s=1, no editing.
    \item If the claim is completely refuted by evidence, i.e. r=1, make editing.
    \item If s<1 and r>0, check whether supported partition is the same as the refuted partition, i.e., check if evidence contradicts with each other.
    \begin{itemize}
        \item If no, they support and refute different partition of the claim, edit the refuted partition.
        \item If yes, then if r>s, this means that with consideration of source reliability and the number of evidence falling into each stance, voice of ''refute'' is stronger than ``support'', we should make editing, otherwise remain original text.
        No matter r>s or s>r in this situation, where there exist controversial reports, we should provide users with two-side evidence over controversial spans.
    \end{itemize}
    \item If s<1 and r=0, this indicates that there are neutral evidence, no editing, but provide users with extra information of this claim.
\end{itemize}
\paragraph{Discussion}
Which factuality label should we use to determine correction?

How should we adjust the criteria to make it more appropriate according to your annotation experience?



\subsection{Edit}
In correction, we keep the principle of making minimal edits against the original text to correct factual errors. 
We do not add extra information provided by evidence that is not directly targeted at factual errors.
No extra deletion, insertion and addition.


\subsection{Merge and deduplicate}
Merge all statements either revised or original ones in order, and deduplicate repeated information with principle of minimal edits.


\section{Annotation}
Based on the annotation guideline, to improve annotation quality and efficiency, we plan to incorporate automatic systems to assist annotation, particularly the subtask of retrieving and collecting evidence for a claim.

An automatic retrieval system is more effective and efficient to collect comprehensive evidence with higher recall than humans, since it takes much less time to go through a large number of documents and extract semantically-relevant snippets of text from a document than an individual annotator. 
So we provide annotators with a set of evidence ranked by relevance degree output by a retrieval system.
Annotators then confirm whether each piece of evidence is closely-related to the claim, if it is, judge the stance, otherwise discard this piece of evidence.
This significantly reduces workload of annotators and improves data collection efficiency.

For other subtasks, such as decomposition, identification of check-worthiness and evidence stance, determination of factual labels and whether correctness is needed, we have tools or models to help produce intermediate results.
It saves time and gets rid of potential inconsistency and calculation errors stemming from individual annotators, at the cost of possibly introducing system bias.

\paragraph{Discussion}
What existing tools, systems, or models can be used for each subtask?
\\
How many examples do we plan to annotate?
\\
In terms of baseline methods, we prefer zero-shot and few-shot learning, or supervised learning?
We set baseline in this work or we need to compete with other approaches?
\\
What is our specific plan for next step?




% Entries for the entire Anthology, followed by custom entries
\bibliography{anthology,custom,ref}
\bibliographystyle{acl_natbib}

% \appendix
% \section{Appendix}
\label{sec:appendix}

\subsection{Answers of ChatGPT}
Let us ask ChatGPT to answer two questions.
We find the answer of the first question makes sense, while the second is nonsensical.

\textbf{when large language models, like ChatGPT tend to hallucinate?}

As a language model, ChatGPT generates text by predicting the most likely next word or sequence of words based on the input it receives. While ChatGPT is trained to generate coherent and grammatical text, there are instances where it may produce text that appears to be "hallucinated" or unrealistic.

The phenomenon of hallucination in language models refers to instances where the model generates text that is either irrelevant, incorrect, or nonsensical in the context of the input it received. This can occur in several situations, including:
\begin{itemize}
    \item Lack of context: ChatGPT relies heavily on context to generate text that is coherent and relevant. If the input it receives is insufficient or ambiguous, it may generate text that is irrelevant or nonsensical.

    \item Over-reliance on training data: ChatGPT is trained on a large corpus of text, which means that it may reproduce text that it has seen before, even if it does not fit the context of the input it received. This can result in text that appears to be "hallucinated" or irrelevant.
    
    \item Insufficient training: If a language model like ChatGPT is not trained on a diverse enough dataset, it may not have the necessary knowledge to generate relevant and coherent text in certain contexts.
    
    \item Bias in the training data: If the training data is biased in some way, such as containing certain stereotypes or patterns, ChatGPT may reproduce these biases in its generated text, leading to unrealistic or biased text.
\end{itemize}
Overall, hallucination in language models is a complex phenomenon that can arise from a variety of factors. While language models like ChatGPT have made significant progress in generating coherent and relevant text, there is still much research to be done to address the issue of hallucination and improve the accuracy of language generation.

\textbf{How to collect questions or generations with objective statements presenting long-tail knowledge?}
To collect questions or generate objective statements that present long-tail knowledge, you can follow these steps:
\begin{itemize}
    \item Identify the topic: Choose a specific topic or niche that you want to explore. This can be anything from a specialized field to a particular hobby.

    \item Conduct research: Use a variety of sources to research the topic. This can include academic journals, online forums, social media, and other relevant resources.
    
    \item Identify long-tail knowledge: Look for information that is not commonly known or readily available. This could be a specific technique or process, a little-known fact, or a unique perspective on the topic.
    
    \item Brainstorm questions and objective statements: Based on your research, brainstorm a list of questions and objective statements that present long-tail knowledge. Focus on topics that have not been extensively covered elsewhere.
    
    \item Refine the list: Review and refine the list of questions and objective statements to ensure that they are clear, concise, and relevant to the topic. Remove any duplicates or irrelevant items.
    
    \item Use appropriate tools: Use appropriate tools such as keyword research tools, online forums, or social media analytics to identify common questions and topics related to the long-tail knowledge you have identified.
    
    \item Analyze and iterate: Analyze the results of your research and adjust your list of questions and objective statements accordingly. Keep refining and iterating until you have a comprehensive list of long-tail knowledge questions and objective statements that you can use to generate content or drive further research.
\end{itemize}




% This is an appendix.

\end{document}
